\chapter{影子银行：管道之外的管道}
\chaptermeta{16}

\begin{ledetext}
有一整套系统，干着和银行一模一样的活，却不穿银行的制服。它不违法。它是这台机器的正常器官，而且拆不掉。
\end{ledetext}

\section{常州那笔 4200 万}

常州武进有家做精密减速器的厂子，老板姓夏。2025 年秋天他接到一个人形机器人关节的订单，要扩一条产线，缺 4200 万。

他先去了银行。信贷员很客气，材料也收得认真：厂房是租的，专用设备做抵押折价折得狠，近三年收入复合增长 40\% 多，但净利润薄，2024 年才刚刚转正。

信贷员没说"不行"。信贷员说的是"我们再研究研究"。

这不是他人品有问题。银行放这笔款要占用资本，权重一算，对银行的资本回报不划算。他手上还有八个客户等着。

夏总最后是在一家美元基金旗下的信贷平台拿到钱的。四年期，前两年只付息，年化 9.8\%，附一堆条款：净负债对 EBITDA 不得超过 3.5 倍，每季度交一次表，超了就重新定价。

这笔钱不是任何人的存款变的。它来自一家北欧的养老金、一家新加坡的保险公司，以及几个连夏总名字都没听说过的家族办公室。

\begin{egbox}{算算他多付了多少}
如果银行肯放，以夏总的资质，报价大概在 5.9\%。他实际拿到的是 9.8\%，中间差 390 个基点。

4200 万 × 9.8\% = 411.6 万，这是他一年的利息。\\ 4200 万 × 5.9\% = 247.8 万，这是他本来可以付的。\\ 一年多掏 163.8 万，四年多掏 655 万。

夏总认了。银行给他的选项不是"更便宜的钱"，是"没有钱"。多付的 390 个基点买的是这条产线今年就能开，而不是等他把利润做厚两年。到那时候订单已经给别人了。

\end{egbox}

\section{银行的三件事，谁都能干}

把银行这门生意拆开，核心只有三件：

\textbf{期限转换。}借进来的是短的，活期存款随时能提；放出去的是长的，房贷三十年。\\ \textbf{流动性转换。}负债端随时能变现，资产端一时半会变不成现金。\\ \textbf{信用转换。}吸进来的是被当作零风险的存款，放出去的是有可能收不回的贷款。

任何机构，只要它同时干这三件事，它做的就是银行的生意。它自称什么不重要。

\term{影子银行}{shadow banking}\index{078@影子银行}的定义就是这么来的：在银行监管框架之外完成这三种转换的一整套机构和市场。

它缺的是两样东西。一是不受资本充足率、拨备覆盖率、流动性覆盖率那一整套约束；二是出事的时候没有存款保险，背后也没有央行的贴现窗口接着。

把持牌银行想成有牌照的出租车：能上专用道，出事有保险，但要装计价器、年检、交份子钱。影子银行是网约车刚出现那两年，同样把人从 A 送到 B，价格更灵活、接单更快，撞了车谁赔当时没人说得清。

这个比喻可以往下推一步：监管的思路从来不是"禁止用私家车拉人"，而是"什么时候发牌、装什么表、买哪种险"。而每一轮发牌落地之后，总会冒出一种还没被写进条例的新送人方式。这不是监管无能，这是定义问题：规则只能写已经存在的东西。

\begin{mythbox}{常见误解}
\mvfrow{误解}{"影子银行就是非法金融、地下钱庄那一套。"}
\mvfrow{事实}{它绝大部分是持牌机构在做合规业务。货币市场基金、资产证券化、信托、私募信贷基金，每一样都有监管，都有信息披露，只是不归"银行"那套监管管。名字里的"影子"说的是它站在银行监管这盏探照灯照不到的角度，不是说它在地下室。真正的非法集资是另一回事，第 23 章会单独拆。}
\end{mythbox}

\section{两条腿，砍掉一条它还站着}

影子银行不是被谁发明出来的，是被两个力量长出来的。

\textbf{第一条腿是监管套利。}巴塞尔框架下，一笔对中型企业的普通贷款风险权重按 100\% 算，按 8\% 到 10\% 的资本要求，放 1 亿要拿出八百万到一千万自有资本压在下面。这笔资本压住了就不能再用。挪到表外，通过一只基金、一个资管计划、一个特殊目的载体去做，资本就腾出来了。同一笔生意，同一个借款人，换个位置放，占用的资本能差一个数量级。

\textbf{第二条腿是真实需求。}银行的风控模型是为标准化客户造的：有抵押、有历史、现金流平稳。夏总这种公司不符合模型，可他的订单是真的，设备是真的，招进来的 40 个人也是真的。这笔钱总得有人给。

只要资本监管还在，只要银行的模型还会筛掉一部分真实的融资需求，影子银行就不可能被消灭。这一句是判断，我把理由摆在这儿了：你可以关掉 P2P，可以让理财净值化，可以把信托的通道砍断，但被赶走的资金和被拒绝的需求会在别的地方接上头，换个名字继续。过去二十年，它至少换过五个名字。

\begin{pullquote}{}
监管能改变影子银行的形状，改变不了它的存在。你堵掉一条通道，水会自己找下一条。
\end{pullquote}

\begin{mythbox}{常见误解}
\mvfrow{误解}{"监管收紧就能把影子银行消灭掉。"}
\mvfrow{事实}{2018 年资管新规之后，中国表外理财的规模确实压下来了，非标转标，通道大幅收缩。但企业的融资需求没有消失，它转到了标准化债权、私募基金、供应链金融和产业链的应收账款里。统计口径变了，链条的形状变了，业务本身还在。监管的现实目标从来不是消灭，是让它可见、可计量、有一层资本垫着。}
\end{mythbox}

\section{过去二十年，它换过几个名字}

\begin{flowlist}
  \flowitem{01}{货币市场基金}{对客户的承诺是"和现金一样，1 美元进 1 美元出"，实际干的是买商业票据。2008 年 9 月 16 日，Reserve Primary Fund 因持有雷曼商票，净值跌到 0.97 美元。承诺碎了一角，之后几天全行业被赎回上千亿美元。它从不自称银行，但它是一家没有存款保险的银行。}
  \flowitem{02}{资产证券化 ABS / MBS}{把一堆房贷打包，按偿付顺序切成不同优先级卖给不同风险偏好的人。逻辑本身没错，今天还在用。错在两处：假设各地房价不会同时跌，以及切完之后拿切剩的部分再切一遍。}
  \flowitem{03}{回购市场}{2008 年真正的挤兑不发生在银行网点。它发生在回购市场：对手方突然要求更高的折扣率，或者干脆不再续做隔夜融资。贝尔斯登和雷曼死于隔夜的钱接不上，不是死于储户排队。这一点在中文报道里几乎从没被讲清楚过。}
  \flowitem{04}{银行理财、信托、委外}{2013 到 2017 年的中国版本。理财资金经信托、券商资管、基金子公司层层通道，投向地方融资平台和房地产。刚兑预期让它在老百姓眼里像存款，底层却是非标债权，期限还错配着。}
  \flowitem{05}{P2P}{高峰期平台数超过五千家。本质上就是没有牌照、没有资本金、没有拨备的银行，负债端全是散户，期限短到一个月。它爆得最惨，因为连"机构投资者"这层缓冲都没有。}
\end{flowlist}

五个名字，五种形态，一套机制。每一次的爆点都落在同一个位置：负债端跑得比资产端快。

\bookbreak

\section{2026 年的主角叫私募信贷}

\begin{statrow}{3}
  \statitem{2.3}{万亿美元}{全球私募信贷管理规模（2025 年）}
  \statitem{17}{\%+}{2020 年以来的年化增速，五年翻倍}
  \statitem{30}{\%}{美国银行流向该领域的敞口占银行资本比重}
\end{statrow}

\term{私募信贷}{private credit}\index{056@私募信贷}就是夏总拿到那 4200 万的地方。不公开发行，不在交易所交易，一家基金直接把钱借给一家公司，条款一对一谈。

钱从哪儿来？养老金、保险公司、大学捐赠基金、家族办公室。这几类机构的共同点是负债期限极长。一家养老金要三十年后才开始大规模付钱，它完全承受得起资产五年不能变现，它想要的是可预期的票息，不是每天跳动的报价。

钱到哪儿去？中型、成长快、抵押品少、现金流还不够漂亮的公司，也就是传统银行覆盖不足的那一段。行业敞口里软件业占三成左右，其中投向 AI 相关领域的约 2000 亿美元，占全部资产 9\%。

价格 8\% 到 12\%，比银行贷款高 200 到 400 个基点；二顺位和夹层 12\% 到 16\%。这个价差不是暴利，它对应的是更差的抵押、更靠后的清偿顺序、更长的锁定期。

\section{银行没有退出，银行绕了个弯}

很多人以为私募信贷的兴起意味着银行在这块业务上撤了。银行确实不直接借给夏总了，但银行借给了借给夏总的那个人。

\begin{chainflow}
\chainnode{你的存款}\chainarrow \chainnode{银行}\chainarrow \chainnode{给私募信贷基金的授信}\chainarrow \chainnode{私募信贷基金}\chainarrow \chainnode{夏总的产线}
\end{chainflow}

这个绕弯叫\term{再中介化}{re-intermediation}\index{080@再中介化}。银行给基金提供循环授信、认购优先级份额，间接参与了企业信贷，账面上记的却是"对金融机构的贷款"，风险权重更低，资本占用更少。

规模有多大：美国银行体系流向私募信贷的资金约 5000 到 6000 亿美元，占银行总资产 3\%，占银行资本 30\%。

\begin{egbox}{把这两个比例做一次除法}
3\% 的资产，对应 30\% 的资本。这两个数放在一起，等于说这批银行的资本对资产比大约是 10\%。它也等于说：银行自有资本的三成，压在这一条通道上。

我见过不少人把"占资本 30\%"念成"银行三成资本要没了"。\textbf{这是错的。}敞口不等于损失，要变成损失得先算回收率。这条通道整体损失 10\%，损失额 550 亿美元，占银行资本 3\%；损失 30\%，占资本 9\%。难受，不是灭顶。

该担心的不是这个数字，是它在四年里从几乎为零涨到了这里，以及没有任何外部机构能确认这些底层贷款是不是用同一套假设估的值。

\end{egbox}

\section{保险公司这一环，问题不在关联交易}

另一件报道不多但很要紧的事：另类资管机构这几年大量收购保险公司，尤其是年金业务。动机很直白，保险公司手里是长期、稳定、成本可控的负债，正好是私募信贷最想要的资金。买下来之后，这些钱可以配置到集团内部发起的资产上。

部分私募股权支持的保险公司，关联方资产配置比例达到 20\% 到 40\%。行业平均水平是 7\%。

\begin{barfig}
  \barrow{私募股权支持的保险公司（区间上沿）}{80}{40\%}
  \barrow{私募股权支持的保险公司（区间下沿）}{40}{20\%}
  \barrow{行业平均}{14}{7\%}
\end{barfig}

问题不在关联交易违法。它合规，有披露，有独立董事签字。问题在估值：你自己发起的贷款，卖给你自己控股的保险公司，价格由你自己的模型给。这条链上没有一个愿意和你对着下注的第三方。

市场价格的全部意义，就是有人愿意在另一边掏钱。没有那个人，剩下的只是一个数字。

\section{把它卖给个人的那一刻}

机构资金和这类资产本来是匹配的：长钱配长资产，锁三年就锁三年。

然后有人想到，这么好的收益率，为什么不卖给个人。

于是有了半流动基金、上市的 BDC（\term{商业发展公司}{business development company}\index{046@商业发展公司}），以及包着这些资产的 ETF。到 2025 年，零售端产品规模约 4000 到 5000 亿美元，占私募信贷总规模两成。

底层贷款要三五年才能收回本金。产品对个人承诺的是季度赎回。

这就是全部的风险所在。

\begin{warnbox}{小心：慢资产配了快负债}
半流动基金一般设有赎回上限（Gate），比如每季度最多赎回净资产的 5\%。平时它是一条摆设，因为正常季度申请赎回的人只有一两个百分点。

它在市场平静时永远不会被触发，在市场紧张时一定会被触发。而它被触发的那一天，就是所有人知道"排队开始了"的那一天。

\end{warnbox}

\begin{egbox}{Gate 是怎么把恐慌放大的}
一只规模 86 亿的半流动基金，季度赎回上限是净资产的 5\%，也就是 4.3 亿。

某个季度它收到 10.4 亿的赎回申请，占净资产 12.1\%。基金按比例配售：每个人只能拿回申请金额的 41\%（4.3 ÷ 10.4）。剩下的 59\% 排队等下一个季度。

下一个季度会发生什么？看到"某基金触发赎回上限"这条新闻的人，比上个季度多得多。原本不打算赎的人也会去排队，因为排在后面的人拿不到钱。

这不是有人在恶意做空。这是每个人做出对自己最合理的选择之后，加总起来的结果。挤兑从来都是这么来的。

\end{egbox}

\begin{talkbox}
  \talkline{新闻}{"某上市 BDC 股价较其最近一期资产净值折价 18\%，创上市以来最大折价。"}
  \talkline{翻译}{这只基金说自己手上的贷款值 100 块，愿意在公开市场上买它的人只肯出 82 块。}
  \talkline{那 18 块是什么}{是市场对两件事同时报价：一，它的内部估值模型可能偏乐观；二，你想立刻拿到现金，就得付这个代价。这个折价是眼下最诚实的温度计，因为它是真金白银交易出来的，不是模型算出来的。}
\end{talkbox}

\section{它是不是下一个次贷}

这个问题被问了两年。两边的论据都不是空穴来风，得两边都摆。

\begin{tablewrap}
\begin{xltabular}{\linewidth}{rXX}
\toprule
\bthead{维度} & \bthead{2007 年的次贷} & \bthead{2026 年的私募信贷} \\
\midrule
\textbf{信用扩张环境} & 长期低利率、全球追逐收益 & 低利率十年之后的高利率，同样在追逐收益 \\
\textbf{风险下沉} & 从优质房贷下沉到次级、无收入证明 & 从大型企业下沉到中型、二顺位、夹层 \\
\textbf{底层透明度} & 打包后无法穿透到单笔贷款 & 不公开交易，按季度按模型估值 \\
\textbf{与银行的耦合} & 银行自持大量超优先级份额 & 银行向基金放贷，敞口占银行资本 30\% \\
\textbf{利率上行的冲击} & 浮动利率重置，还款额跳升 & 多为浮动利率，借款人利息负担同步上升 \\
\textbf{嵌套层数} & CDO、CDO 平方，多层再证券化 & 基本没有再证券化 \\
\textbf{杠杆倍数} & CDO 10–20 倍，投行 20–30 倍 & 基金层面 0–2 倍 \\
\textbf{负债端稳定性} & 隔夜回购、ABCP，一夜之间能抽干 & 养老金保险的多年期承诺，锁定期长 \\
\textbf{风险释放速度} & 几周 & 按季度重估，几个季度到几年 \\
\bottomrule
\end{xltabular}
\end{tablewrap}

上半张表说明它像。下半张表说明它不像。两边都是真的。

中金把它称为"2 万亿美元的灰犀牛"，判断是风险可能在未来 12 到 24 个月逐步传导，但目前不构成黑天鹅式的系统性危机。

这句话要按它本来的分量读：\textbf{它是一个判断，不是一个事实。}它站在三个假设上：负债端不会突然变短，零售化比例不会失控，投向 AI 的那 2000 亿不会集体出问题。三个里第一个最结实，第三个最悬。

我的看法：如果这轮真出事，触发点大概率不在企业违约率本身。企业是慢慢坏的，一个季度坏一点，机构消化得了。快的是产品端的赎回。快的一头撞上慢的一头，中间那条缝就是风险。

这是对机制的判断。什么时候发生、会不会发生，我不知道。

\section{你能盯的四个数}

这行有个好处：它的体温计是公开的。

\begin{flowlist}
  \flowitem{01}{BDC 的价格对净值比}{日频，公开可查。折价持续扩大并且伴随成交量放大，说明市场在给模型估值打问号。单日折价没意义，连续两个季度走扩才是信号。}
  \flowitem{02}{高收益债利差和杠杆贷款价格}{私募信贷自己不交易，它的表兄弟天天交易。杠杆贷款平均报价跌破 95 通常意味着压力开始了。表兄弟先咳嗽。}
  \flowitem{03}{大型银行和另类资管机构的 CDS 报价}{市场为"这家机构会不会出事"直接标的价，比任何评级都灵敏，因为报价的人要用自己的钱兑现。}
  \flowitem{04}{赎回上限有没有被触发}{这类消息通常不上头条，藏在产品公告里。第一家触发的时候，市场往往还很平静。}
\end{flowlist}

说清楚：这四个数不是买卖信号，本书不提供投资建议。它们唯一的用处，是让你在新闻标题变得吓人之前，先知道水位到了哪一格。

\section{中国这边是反过来走的}

2018 年资管新规是个分水岭。打破刚兑，产品净值化，非标转标，多层嵌套的通道被砍断。信托被迫转型，从"影子银行的通道"往资产服务信托、资产管理信托、公益慈善信托这三类走，很多公司到今天还没转完。

地方融资平台的债务在化解，手段是展期、置换、降息：短的换成长的，贵的换成便宜的，隐性的挪到显性的账上。这不消灭债务，它买时间。

方向和美国正好相反。美国是银行受资本约束往后退，把资产让给非银机构；中国是先长出一个庞大的影子银行体系，再用行政力量把它往表内和标准化市场里赶。

共同点只有一个：无论赶到哪边，那笔钱最后还是要落到某个像夏总一样的人手上，去买一台设备，雇几十个人。中间过几道手，决定的是谁承担损失，不是钱去干什么。

\begin{takeaway}{本章一句话}
影子银行不是银行的敌人，是银行监管的影子。它干的是同一件事：拿短钱做长事。它省下的是资本占用，代价是出事时没有存款保险，也没有最后贷款人。

\begin{itemize}
  \item 判断标准只有一条：谁在做期限、流动性、信用这三种转换，谁就在做银行的生意。
  \item 它有两条腿：监管套利和真实需求。这两条腿在，它就永远在，只是换名字。
  \item 2026 年它最大的一块叫私募信贷，2.3 万亿美元。风险不在违约率，在慢资产配了快负债。
  \item "占银行资本 30\%"是敞口不是损失。看数据要先弄清楚它是哪一种。
\end{itemize}

\end{takeaway}

\begin{bridgebox}
\textbf{下一章}影子银行自己不会引爆什么。它得先和另一样东西凑到一起：一个所有人都相信、而且确实是真的新故事。这两样一凑，就会重演一部上演了四百年的老剧本。
\end{bridgebox}

